% page 515 du fac-simile (image p-514 du scan)
% bloc 182.9 x 263.6 mm ; marge gauche 19.1 mm
\pgc{-2.540mm}{0.000mm}{
\l{\marge[86.589mm]{34.178mm}{\ornamento{10.817mm}{ornements/letroj/litero-S.png}}}
\l{}
\l{}
\l{}
\l{\cel{63}507}
\l{}
\l{}
\l{}
\l{}
\l{}
\l{}
\l{}
\l{\cel{4}\sou{sabato}. (\sou{biblo}).Repozo, preskriptita da lia religio, quan la Judi}
\l{\cel{7}devas praktikar dum la tota dio sepesma di la semano. - DEFIRS.}
\l{}
\l{\cel{4}\sou{sabeliko}. (\sou{bot.)}\cel{3}Varietato di kaulo, di qua la folii esas tre}
\l{\cel{7}frizita. L.\cel{1}\sou{brassica oleracea Sabauda}.- I.}
\l{}
\l{\cel{4}\sou{sablo}.- Substanco pulveratra, produktita da la desagregeso di la}
\l{\cel{7}roki silicoza o quarcoza. - FIS.}
\l{}
\l{\cel{4}\sou{sabordo}. (\sou{navo)}. Aperturo quadri-angula, situita en la kareno di}
\l{\cel{7}navo (specale por lasar pasar la boko di kanono). - F.}
\l{}
\l{\cel{4}\sou{sabro}. - Armo ne-pafala, kun pinto ed aristo tranchiva, la aristo}
\l{\cel{7}tranchiva di la lamo esante konvexa. - DEFIRS.}
\l{}
\l{\cel{4}\sou{sacerdoto}. I. Ta qua prezidas la kulto deala, la ceremonii reli-}
\l{\cel{7}giala. - II. (\sou{religio katol.}) Ta persono qua recevis de la epis-}
\l{\cel{7}kopo, kun la ordo di sacerdoteso, la permiso celebrar la meso e}
\l{\cel{7}donar la uncioni. - EFIS.}
\l{}
\l{\cel{4}\sou{saciar.}\cel{2}(\sou{trans.}) Kontentigar per cesigar la bezono. - FIS.}
\l{}
\l{\cel{4}\sou{sadismo}. Laciveso akompanata da krueleso. - DEFIRS.}
\l{}
\l{\cel{4}\sou{safena}. (\sou{anat.})\cel{2}Veino di la gambo, qua departas de la pedo-fingri.}
\l{}
\l{\cel{4}\sou{safiro}. Lapido briloza, bele blua.- DEFIRS.}
\l{}
\l{\cel{4}\sou{safrano}. I. (\sou{bot}.)Planto bulboza di qua la floro, blua kun purpurajo}
\l{\cel{7}havas stigmati flavatra qui, pulverigita, uzesas kom kolorizivo}
\l{\cel{7}e kom kondimento. - L.\cel{1}\sou{safranum}. - II. La ipsa farbo. - DEFIRS.}
\l{}
\l{\cel{4}\sou{sagaca}. Karakterizata da instinto subtila pri deskovrar la kozi,}
\l{\cel{7}pri penetrar til la fakti, qua komprenas la relati, mem minima,}
\l{\cel{7}inter la kozi ipse, qua dicernas multa diferi ube altra personi}
\l{\cel{7}vidas nur kozi uniforma.- EFIS.}
\l{}
\l{\cel{4}\sou{sagitario}. (\sou{bot.)}\cel{2}Planto qua vivas en aquo, herbatra, perena,}
\l{\cel{7}ula folii di qua esas flecho-forma e vertikala, la ceteri longa}
\l{\cel{7}e flotacanta. - L.\cel{1}\sou{sagittaria}. - FISL.}
\l{}
\l{\cel{4}\sou{saguino}. (\sou{zool.)}\cel{3}Speco di mikra simio, kun kaudo longa. -}
\l{\cel{7}L.\cel{1}\sou{rapale}. - DEFIRS.}
\l{}
\l{\cel{4}\sou{saguto.}\cel{1}(\sou{bot.)}\cel{1}Palmiero di qua la medulo donas fekulo amilatra,}
\l{\cel{7}e la frukti donas, per distilo, liquoro ebriigiva.-\cel{2}DEFIRS.}
\l{}
\l{\cel{4}saimo. Graso de la epiplono di la porko, fuzita e salizita, quan}
\l{\cel{7}onu uzas koquarte, por fritar, e c. - I.}
}
